\section{Розділ~\ref{sec:sensors}. Входи машини}

Будову датчиків виміряно прямо --- записами струму поодиноких клітин, коливань завитки й підрахунком клітин сітківки; межа чутливості й формула дробового шуму випливають із фізики, а те, що коди датчиків підлаштовано під найбільшу інформацію, --- гіпотеза з сильною опорою.

{\footnotesize
\setlength{\tabcolsep}{3pt}\renewcommand{\arraystretch}{1.15}
\begin{longtable}{>{\raggedright}p{0.5cm}>{\raggedright}p{4.0cm}>{\raggedright}p{5.6cm}>{\raggedright}p{2.3cm}>{\raggedright\arraybackslash}p{1.7cm}}
\toprule
№ & Твердження & Дослід і що виміряно & Джерело & Статус \\
\midrule\endfirsthead
\toprule
№ & Твердження & Дослід і що виміряно & Джерело & Статус \\
\midrule\endhead
1 & Датчик --- перетворювач: рецептор дає струм, а вихід чутливих клітин ока й вуха --- глутамат на шину \nt{GLUT}; перші клітини імпульсів не видають (рис.~\ref{fig:transducer}) & Палички й колбочки черепахи виділяють збуджувальну амінокислоту: її вловлюють глутаматні канали на відокремленому клаптику мембрани. Внутрішні волоскові клітини завитки щура: струми в закінченнях нервових волокон ідуть через глутаматні AMPA-рецептори, частота вивільнень зростає з плавною деполяризацією клітини до $150$\,Гц & \cite{CopenhagenJahr1989,GlowatzkiFuchs2002} & виміряно \\
2 & Датчик світла в темряві відповідає на один фотон струмом близько пікоампера & Всмоктувальний електрод на зовнішньому сегменті палички ропухи: відповіді на тьмяні спалахи квантовані, подія $\approx1$\,пА ($5\%$ темнового струму) з розкидом $0.2$\,пА, ефективність $0.5$. Людина повідомляє про влучання одного фотона частіше, ніж випадково & \cite{Baylor1979,Tinsley2016} & виміряно \\
3 & Датчик звуку помічає зміщення менше за нанометр & Метод Мессбауера, основна мембрана завитки морської свинки в місці $16$--$19$\,кГц: на порозі відповіді слухового нерва швидкість мембрани близько $0.04$\,мм/с, тобто зміщення $\approx0.35$\,нм (наш перерахунок). Вимірювали мембрану, а не пучок датчика & \cite{Sellick1982,Hudspeth1989} & виміряно \\
4 & Точність у темряві обмежена дробовим шумом фотонів~\eqref{eq:photonnoise}; за слабкого світла накопичення довше & Формула випливає з пуассонівської статистики. У паличці шум за тьмяного світла збігається із сумою незалежних квантових подій; на яскравому тлі відповідь на спалах менша й коротша. Число «десяті частки секунди» тут окремо не перевірено & \cite{Baylor1979} & теорія \\
5 & Діапазон входу --- десять порядків для світла й дванадцять для звуку, а частота імпульсів --- менше за три порядки & Для звуку: відповідь основної мембрани на характеристичній частоті зростає з крутістю до $0.2$\,дБ/дБ у смузі $40$--$80$\,дБ, стиснення видно й вище $100$\,дБ. Самі числа порядків --- стандартні оцінки, у розділі без джерела & \cite{Ruggero1997} & виміряно \\
6 & Відповідь датчика~\eqref{eq:nakarushton}, а $\sigma$ іде за середньою яскравістю, зсуваючи криву логарифмічною віссю (рис.~\ref{fig:adaptation}) & Формулу вперше підібрано до S-потенціалів сітківки риби. Колбочки черепахи: відповіді підкоряються $V=I V_m/(I+\sigma)$ і охоплюють $\sim3.5$ порядку яскравості; після $2$--$3$ хвилин тла крива зсувається по горизонталі, зберігаючи ширину. Зв'язок із діленням на загальну активність --- огляд & \cite{NakaRushton1966,NormannPerlman1979,Carandini2012} & виміряно \\
7 & Датчики передають зміни, а їхні коди підлаштовано під найбільшу інформацію на імпульс (твердження~\ref{prop:sensors}) & Принцип запропоновано теоретично. Найсильніше свідчення: у мухи характеристика «контраст $\to$ відповідь» нейронів першого шару збігається з кумулятивним розподілом контрастів природних сцен --- так досягається найбільша інформація & \cite{Barlow1961,Laughlin1981} & гіпотеза \\
8 & Датчики ввімкнення й вимкнення --- двопровідний код; подієва камера повторює його в кремнії & Огляд дослідів: канали ввімкнення й вимкнення сітківки розділяються вже на першому синапсі й вимикаються окремо. Камера: $128\times128$ пікселів без кадрів, діапазон $120$\,дБ, затримка $15$\,мкс & \cite{Schiller1992,Lichtsteiner2008,Gallego2022} & виміряно \\
9 & В оці $\sim10^8$ датчиків світла й $\sim10^6$ вихідних нейронів: стиснення в $\sim100$ разів & Підрахунок на цілих препаратах сітківки людини: у середньому $92$ млн паличок і $4.6$ млн колбочок; вихідних (гангліозних) клітин від $0.7$ до $1.5$ млн у різних очах & \cite{Curcio1990,CurcioAllen1990} & виміряно \\
10 & У миші вихідних каналів ока понад тридцять & Миша: двофотонна кальцієва зйомка понад $11\,000$ вихідних клітин і кластеризація відповідей --- помітно більше за $30$ функціональних типів. Для людини число не виміряно & \cite{Baden2016} & виміряно \\
11 & Око морської свинки передає $\sim875\,000$ біт/с; перерахунок на око людини дає $\sim10^7$ біт/с & Морська свинка, багатоелектродна матриця, рух у природних сценах: $6$--$13$ біт/с на клітину, $\sim875\,000$ біт/с на $\sim10^5$ вихідних клітин. Для людини ($\sim10^6$ клітин) $\sim10^7$ біт/с --- перерахунок авторів & \cite{Koch2006} & виміряно \\
12 & Вухо --- банк смугових фільтрів із майже логарифмічною картою частот (рис.~\ref{fig:filterbank}) & Місце найбільшого коливання основної мембрани залежить від частоти; функція «частота--місце» майже експоненційна й однією формою описує завитки людини та живих тварин & \cite{Greenwood1990,Robles2001} & виміряно \\
13 & Резонатори активні: частина датчиків підсилює слабкі коливання в тисячі разів за амплітудою ($66$--$76$\,дБ), зі зростанням гучності підсилення падає & Лазерна віброметрія основи завитки шиншили: за слабких тонів на характеристичній частоті підсилення $66$--$76$\,дБ відносно стремінця; смерть знижує чутливість на $60$--$81$\,дБ (у $10^3$--$10^4$ разів за амплітудою) і прибирає стиснення. Джерело сили --- зовнішні волоскові клітини & \cite{Ruggero1997,Robles2001} & виміряно \\
14 & Підсилювач вуха стоїть біля межі самозбудження & Розрахунок: біля точки біфуркації Хопфа неминучі стиснення діапазону, гостре налаштування й комбінаційні тони --- усі відомі нелінійності слуху. Що завитка налаштована саме так, прямо не доведено & \cite{Eguiluz2000} & гіпотеза \\
15 & На низьких частотах імпульси йдуть у фазі зі звуком (у морської свинки прив'язка слабшає від $\sim600$\,Гц і помітна до $\sim3.5$\,кГц), і за ними знаходять напрямок; на високих лишається код місцем & Слуховий нерв морської свинки: прив'язка до фази починає падати близько $600$\,Гц і зникає вище $3.5$\,кГц; у кішки й мавпи межа приблизно на октаву вища. Різниця часу приходу у два вуха --- огляд & \cite{PalmerRussell1986,Grothe2010} & виміряно \\
16 & Інші датчики (табл.~\ref{tab:sensors}): у нюху $\sim400$ типів рецепторів, по одному на нейрон, і комбінаторний код; смак частково через АТФ і серотонін; дотик --- датчики, що звикають швидко й повільно; тепло й холод --- окремі & Миша: один рецептор упізнає багато запахів, один запах --- багато рецепторів, різні запахи --- різні набори. Без рецепторів АТФ нерв смаку не відповідає; смакові бруньки виділяють серотонін. Запис поодиноких волокон шкіри руки людини: чотири типи за швидкістю звикання. Холодовий канал відмінний від теплових & \cite{Mombaerts2004,Malnic1999,Finger2005,Huang2005,JohanssonVallbo1979,McKemy2002} & виміряно \\
\bottomrule
\end{longtable}}
